El radó 222, de símbol Rn, és un gas noble responsable de bona part de l'exposició de les persones a les radiacions ionitzants. El $^{222}\mathrm{Rn}$ es forma al subsòl a partir del radi (Ra) i a causa del seu estat gasós es difon cap a l'atmosfera.

\begin{apartat}{1}
Quan el $^{222}\mathrm{Rn}$ es desintegra emet partícules $\alpha$. Escriviu l'equació nuclear d'aquest procés de desintegració.
\end{apartat}

\begin{apartat}{1}
A més de la radiació $\alpha$, durant el procés de desintegració també s'emeten raigs $\gamma$ (no cal que els inclogueu en l'equació de l'apartat anterior). Calculeu la freqüència i la longitud d'ona d'un fotó $\gamma$ d'energia $5,50\ \mathrm{MeV}$.
\end{apartat}

\dades{%
  Nombres atòmics: Bi, 83; Po, 84; At, 85; Rn, 86; Fr, 87; Ra, 88; Ac, 89.\\
  $1\ \mathrm{eV} = 1,60\times 10^{-19}\ \mathrm{J}$\\
  Constant de Planck, $h = 6,63\times 10^{-34}\ \mathrm{J\,s}$\\
  Velocitat de la llum, $c = 3,00\times 10^{8}\ \mathrm{m\,s^{-1}}$}
